\section{Pre-action prompt repair and exact error accounting}
\label{sec:formulation}

% Methods source map: M01. EXP-016, EXP-034--037, EXP-056.
\subsection{A finite decision with an observation boundary}
A state $h$ identifies an image--class pair, its prompt history, and already
computed masks. SAM~3.1 remains frozen. Each executed base action returns a
set of retained candidate masks: $\mathcal C_0(h)$ for canonical text,
$\mathcal C_1(h)$ for normalized text, and $\mathcal C_2(h)$ for the fixed
text template. A0 and A1 use identical text in this vocabulary but retain
separate registered call identities; equality of their returned masks is
not assumed. Thus three base calls do not imply three distinct semantic
prompts. We use $A_i$ for the action and its union-mask outcome when
the context is unambiguous. The remaining vocabulary is
\begin{equation}
 \mathcal A(h)\subseteq\{\operatorname{STOP},A3,A4,A5,A6\},
 \label{eq:unified-action-set}
\end{equation}
with structural feasibility determined by the frozen action plan. These
repairs apply a horizontal flip, a source-candidate crop, a source-candidate
box, or a source-candidate point, respectively. STOP retains the pixelwise
union of \emph{all} retained A0 candidates and requires no new call; it does
not select the highest-scoring candidate. The latter is used only to plan
source-dependent repairs. Appendix~\ref{sec:implementation-details} specifies
the prompts, restoration rules and distinction between structural
infeasibility, valid empty outcomes and technical missingness.

Let $\mathcal F_0(h)$ denote the information generated by registered metadata,
the executed base actions, and their observable masks and status. A prediction
used to choose an unexecuted action must be measurable with respect to this
information. Reference labels, future masks, and a candidate action's realized
runtime may define training targets or evaluation records, but are excluded
from the prospective predictor. This is an observation-time requirement, not
an independence assumption between observed masks and reference labels.
Pre-GT label blindness alone does not establish this action-time
observability; the offline failure-handling exception is described in
Section~\ref{sec:missingness-timing}.

Write $R(h,a)$ for the realized residual and $c_{\rm pre}(a)$ for the fixed
training-derived action-type cost profile. At a registered multiplier
$\lambda$, the value and selected action are
\begin{equation}
 V_\lambda(h,a)=R(h,a)+\lambda c_{\rm pre}(a),\qquad
 \widehat a_\lambda(h)=\mathop{\arg\min}_{a\in\mathcal A(h)}
       \{\widehat R(h,a)+\lambda c_{\rm pre}(a)\}.
 \label{eq:unified-value}
\end{equation}
The oracle $a^*_\lambda(h)$ minimizes the corresponding realized value using
the same feasible actions and action order. It is an evaluation reference
that consumes action outcomes, not an available pre-action policy.

% Methods source map: M02. EXP-016 and prospective protocol boundary.
\subsection{Observable geometry and target-only validity}
\label{sec:observable-geometry}
The representation must describe spatial variation without seeing a future
repair (Fig.~\ref{fig:representations}). A scalar error target does not locate that variation. We therefore
use the candidates already available at the decision boundary to define
local supports. Their agreement pattern describes observable mask membership;
it does not label semantic objects or identify causal effects. Connectivity
then separates disconnected occurrences of the same pattern, allowing a
learner to associate error with local area, position and shape.

Let $\Omega_h$ be the full image raster and let
$(M_1,\ldots,M_{J_h})$ enumerate every retained candidate in
$\mathcal C_0(h)\cup\mathcal C_1(h)\cup\mathcal C_2(h)$, with candidate
identity preserved and a deterministic ordering. Each pixel has signature
$s_h(x)=(\mathbf 1\{x\in M_j\})_{j=1}^{J_h}$. Thus $J_h$ varies with the
candidate count; these are not three bits formed from the three action
unions. Splitting each signature class into maximal four-connected
components produces a disjoint partition
$\mathcal B(h)=\{B_k(h)\}_{k=1}^{K_h}$ of $\Omega_h$. Full-raster areas,
positions, shapes, and agreement descriptors are observable. Deterministic
component identifiers and their label-map indices must refer to the same
pixel supports. The all-zero signature retains the complement of the
candidate union. With no retained base candidates the entire raster is one
cell. At the F1 observation, retained A3 candidates extend this same
candidate-level construction.

Evaluation may additionally declare a target-valid subset
$G_h\subseteq\Omega_h$. For example, VOC void pixels are excluded from the
target. This subset must not alter prospective geometry or prediction
weights. With reference label $Y_h$ and realized action mask
$\widehat Y_{h,a}$, define
\begin{align}
 e_{h,a}(x)&=\mathbf 1\{\widehat Y_{h,a}(x)\ne Y_h(x)\},\nonumber\\
 R(h,a)&=\frac{\sum_{x\in G_h} e_{h,a}(x)}{|G_h|},\qquad
 r_k(h,a)=\frac{\sum_{x\in B_k(h)\cap G_h}e_{h,a}(x)}{|B_k(h)\cap G_h|}.
 \label{eq:unified-targets}
\end{align}
The local target is undefined when its denominator is zero. An unexecuted
candidate mask need not be constant on an observable cell: $r_k$ is its
mean mismatch over that cell, not a binary cell label.

\begin{proposition}[Exact accounting on a fixed partition]
\label{prop:unified-accounting}
For a nonempty valid support, summing local mismatch counts over the disjoint
cells exactly reconstructs the direct mismatch count. Consequently,
\begin{equation}
 R(h,a)=\sum_{k:\,|B_k\cap G_h|>0}
       \frac{|B_k\cap G_h|}{|G_h|}\,r_k(h,a).
 \label{eq:unified-exact-accounting}
\end{equation}
\end{proposition}
\begin{proof}
The sets $B_k\cap G_h$ partition $G_h$. Therefore
$\sum_{x\in G_h}e_{h,a}(x)=\sum_k\sum_{x\in B_k\cap G_h}e_{h,a}(x)$.
Substitution of the nonempty-cell means and division by $|G_h|$ give the
identity. Empty intersections contribute zero counts.
\end{proof}

The proposition applies to any disjoint exhaustive partition; it does not
establish a unique benefit of Boolean atoms. It also does not grant access to
$G_h$ at prediction time. The prospective local model instead aggregates as
\begin{equation}
 \widehat R(h,a)=\sum_k\alpha_k(h)\widehat r_k(h,a),\qquad
 \alpha_k(h)=|B_k(h)|/|\Omega_h|.
 \label{eq:unified-prospective-aggregation}
\end{equation}
These full-raster weights are distinct from the target-valid weights in
Equation~\ref{eq:unified-exact-accounting}. Exact target accounting therefore
does not imply exact learned prediction. The corrected bank audit, including
its integer-count check, is reported separately from learned-model evidence
in Supplement~\ref{sec:accounting-appendix}.
The learned local outputs are an intermediate representation, not a separate
demonstration of local calibration. In particular, prospective first-batch R1\_P uses only
the state residual for supervision. The region-count-matched rectangle and
candidate-union controls test alternatives to Boolean localization; their
feature content also differs, so they do not isolate boundaries alone.

% Methods source map: M03. EXP-034--036, EXP-038, EXP-056--057.


\subsection{Absolute fidelity, relative fidelity, and utility}
Mean absolute error (MAE) measures errors in $\widehat R$. For a defined non-STOP action
and STOP anchor, define
\begin{equation}
 \Delta(h,a)=R(h,a)-R(h,\operatorname{STOP}),\qquad
 d(h,a)=|\widehat\Delta(h,a)-\Delta(h,a)|.
 \label{eq:unified-drre}
\end{equation}
The mean STOP-relative error is termed decision-relative residual error
(DRRE). Let $\mathcal D$ collect the phase's defined state--action cells, and
$\mathcal P$ its defined non-STOP cells with a defined STOP anchor. The
prospective metrics are pooled record means,
\begin{equation}
 \operatorname{MAE}=\frac{1}{|\mathcal D|}
     \sum_{(h,a)\in\mathcal D}|\widehat R(h,a)-R(h,a)|,\qquad
 \operatorname{DRRE}=\frac{1}{|\mathcal P|}
     \sum_{(h,a)\in\mathcal P}d(h,a).
 \label{eq:unified-fidelity-aggregation}
\end{equation}
V6 uses its complete-case state rule; COCO retains defined fidelity cells
under its frozen technical-missingness rule. Neither is a mean of per-image
ratios. Image groups determine the paired resampling unit. The weights used
to train a decision-relative loss are not additional metric weights.
Equivalently, the pooled DRRE weights each state's within-state mean by
its number of defined non-STOP pairs. It is generally different from an
equally weighted mean of state means when those counts vary; the
state-balanced training loss does not change this evaluation definition.
The historical diagnostic manuscript defined an equally weighted mean of
state means, rather than this pooled expression. Its actual execution
reduction remains unverified. We therefore keep those historical results
separate from the prospective comparisons under Equation~\ref{eq:unified-fidelity-aggregation}.
A common additive prediction error
across actions cancels in $\widehat\Delta$; a small absolute error can still
reverse a close ordering. This motivates measuring both quantities without
claiming either is a universal proxy for policy value.

For the frozen complete-oracle state population $\mathcal H$, normalized
regret is the ratio of aggregated regret to aggregated STOP regret:
\begin{equation}
 \operatorname{NR}_\lambda=
 \frac{\sum_{h\in\mathcal H}[V_\lambda(h,\widehat a_\lambda(h))-
                   V_\lambda(h,a^*_\lambda(h))]}
      {\sum_{h\in\mathcal H}[V_\lambda(h,\operatorname{STOP})-
                   V_\lambda(h,a^*_\lambda(h))]}.
 \label{eq:unified-normalized-regret}
\end{equation}
Positive oracle headroom makes STOP the unit reference. The COCO E4
\emph{utility improvement} is exactly $1-\operatorname{NR}_0$. Its interval
concerns this dimensionless normalized quantity, not raw realized gain or
IoU improvement. A zero denominator makes the metric undefined.

Sparse dispatch ranks states by predicted STOP value minus the best feasible
non-STOP value. For an eligible population of size $N>0$ and budget $b$,
the primary forced rule selects
$K=\max\{1,\lceil bN\rceil\}$ eligible states, using the frozen ordering and
state-ID tie break, and executes their best predicted non-STOP actions even
when predicted gains are nonpositive. If $S_K$ is that selected set, then
\begin{equation}
 \operatorname{GC}(b)=
 \frac{\sum_{h\in S_K}[V_\lambda(h,\operatorname{STOP})-
                          V_\lambda(h,\widehat a^+_\lambda(h))]}
      {\sum_{j=1}^{K}g^*_{(j)}},\qquad
 g^*(h)=V_\lambda(h,\operatorname{STOP})-V_\lambda(h,a^*_\lambda(h)).
 \label{eq:unified-gain-capture}
\end{equation}
The oracle gains are nonnegative because STOP is feasible, and are ordered
from largest to smallest on the identical eligible population. Negative
realized gains remain in the numerator. The secondary positive-gain cap may
use fewer states, while retaining the same budget-specific oracle denominator.
For COCO it has no role in criterion E4-D, E4 support, or E5 classification.

\subsection{Separating state selection from action choice}
\label{sec:selector-action-diagnosis}
Sparse repair contains two decisions. A selector allocates the limited calls
across states; an action head chooses a repair within each selected state.
Accurate action ordering at an unselected state produces no gain, while
selecting a promising state can still lead to the wrong repair.
Table~\ref{tab:selector_action_framework} defines four oracle-substitution
diagnostics that separate state allocation from within-state action choice.
Any numerical use must hold the eligible population, feasible actions,
budget semantics and oracle denominator fixed. The table defines the
diagnostic questions; it does not report a newly evaluated decomposition.

Four elementary conditions clarify a complementary, within-state question
(Supplement~\ref{sec:explanatory-propositions}). A common prediction offset
preserves relative values and the action choice
(Proposition~\ref{prop:common-mode-invariance}). A strict bound on an action's
STOP-relative error preserves its beneficial sign
(Proposition~\ref{prop:stop-sign-stability}). With a unique optimum, a margin
larger than twice a uniform action-value error bound preserves the entire
choice (Proposition~\ref{prop:action-identity}); the same bound limits regret
to decisions with small margins (Proposition~\ref{prop:margin-regret}).
These explanatory statements do not supply the required error
bound for a fitted model. Small mean MAE or DRRE is insufficient, and
within-state stability does not guarantee across-state sparse ranking.
The historical VH3 margin mechanism did not replicate
(Supplement~\ref{sec:mechanisms-history}).
